\documentclass{article}

% ============================================================
% INFR11269 Machine Learning Systems — Coursework Report
% Academic year 2026/27
%
% This coursework uses the official ICLR 2027 LaTeX style.
%
% IMPORTANT FORMAT RULES
% - Do not edit iclr2027_conference.sty.
% - Do not change page geometry, text width/height, base font size,
%   or line spacing.
% - Do not enable \iclrfinalcopy.
% - Normal content packages are fine if they do not alter the report format.
%
% Main-body page limit: six pages.
% References and Appendix do NOT count toward the main-body page limit.
% Reviewers are not obliged to read Appendix at all. Please put your main content
% in the main body.
% 
% AI policy:
% Academic integrity is an underlying principle of research and academic practice.
% All submitted work is expected to be your own, including any material you have
% generated yourself with the assistance of AI tools such as ELM, ChatGPT,
% Bing Chat, Claude, Google Gemini, or any other automated assistance.
% For all such material you must make prominent explicit acknowledgement and
% explain clearly what you have done in the AI declaration section.
% This includes using AI to directly generate text, code, figures or similar
% or to complete other tasks. In all cases, you remain entirely responsible
% for the material you produce.
% ============================================================

\usepackage{iclr2027_conference,times}

\usepackage{hyperref}
\usepackage{url}
\usepackage{booktabs}
\usepackage{graphicx}
\usepackage{amsmath}
\usepackage{amssymb}

% The official ICLR archive includes math_commands.tex. It is optional:
\input{math_commands.tex}

\title{Custom CUDA Kernels for GLM-ASR and RAG System Investigation}

% Keep this line unchanged.
\author{Anonymous coursework submission}

% Keep this COMMENTED OUT.
% Review mode is part of the coursework format and is checked automatically.
% \iclrfinalcopy

\begin{document}

\maketitle

% Course-specific template fingerprint. Keep this line unchanged.
\lhead{Under review as a report for UoE 26/27 Sem1 INFR11269: Machine Learning Systems}

% Keep this COMMENTED OUT.
% \begin{abstract}
% \end{abstract}

% ============================================================
% REPORT CONTENT GUIDANCE

% The structure below is RECOMMENDED, not mandatory.
% You may rename, merge, split, or reorder sections if another structure
% better communicates your investigation.

% Regardless of section titles, the report should make it easy to find:
%   1. the problem, motivation, and relevant system context;
%   2. the supplied baseline and your substantive technical work;
%   3. experimental methodology, metrics, and controlled comparisons;
%   4. results covering task/system quality AND relevant system cost;
%   5. analysis of why the results occur, limitations, and conclusions.
% Well organized report would result in better scores.

% Depth is preferred over implementing many unrelated features.
% ============================================================

\section{Introduction}
\label{sec:introduction}

This report covers INFR11269 coursework on machine learning \emph{systems}:
implementing GPU kernels for GLM-ASR-Nano speech recognition and investigating a
retrieval-augmented generation (RAG)~\cite{lewis2020rag} pipeline built on that stack. The ASR model
routes heavy operations through four CUDA tasks while PyTorch manages tensors and
orchestration. For RAG, we study the question: \emph{under a limited context
budget, how does evidence selection affect answer quality and context/cost?}

The ASR baseline is the supplied skeleton with empty blanks and PyTorch fallbacks
(\texttt{./run.sh --cuda none}). We implemented all four kernel groups, verified
each with \texttt{make test1}--\texttt{make test4}, and checked the full model with
\texttt{./run.sh} (all CUDA groups). For RAG, we established a practice-split
retrieval baseline (350 queries), ran a controlled practice pilot, and completed a
held-out \textbf{dev validation} (125 queries, five evidence-selection
configurations) without tuning on dev after observing scores.

To satisfy the ASR reporting requirement, one selected question answer per task is
included in Section~\ref{sec:results}: Task~1 in
\textit{Selected ASR question answer (Task 1)}, and Tasks~2--4 in
\textit{Selected ASR question answers (Tasks 2--4)}. Optional additional Task~1
analysis is reported separately. Reproduction commands and additional records are
in the Appendix.

\section{System and Technical Approach}
\label{sec:approach}

\subsection{Supplied baseline}
The homework provides \texttt{csrc/task*.cu} with numbered blanks,
\texttt{launch\_*} grid configuration, CPU references in \texttt{tests/}, and
\texttt{glm\_asr\_cuda/ops.py} which selects CUDA or PyTorch per kernel group.

\subsection{Task 1: RoPE kernel design}
RoPE rotates pairs \((v[i], v[i+h])\) before attention, with $h=\mathrm{rotary\_dim}/2$:
$\angle = \mathrm{pos}\cdot \mathrm{inv\_freq}[i]$,
$v[i]' = a\cos\angle - b\sin\angle$,
$v[i+h]' = b\cos\angle + a\sin\angle$.
One launch covers both \texttt{q} and \texttt{k}: threads $0,\ldots,\text{q\_pairs}-1$
touch \texttt{q}, then \texttt{k}. Each thread maps a global pair index to
\texttt{row}, pair index \texttt{i}, and token position
$\mathrm{pos} = (\mathrm{row}/\mathrm{num\_heads}) \bmod \mathrm{seq\_len}$.
Elements beyond \texttt{rotary\_dim} are untouched (audio encoder uses
\texttt{rotary\_dim}$=32$ of \texttt{head\_dim}$=64$).

Implementation choices: \texttt{ROPE\_THREADS}$=256$ blocks sized by
\texttt{ceil(total\_pairs / 256)}; in-place updates; \texttt{cosf}/\texttt{sinf}
computed per pair (no lookup table). Blanks~1.1--1.4 implement indexing and the
rotation; \texttt{launch\_rope} is unchanged.

\subsection{Task 2: warp/block reductions, RMSNorm, softmax}
Part~A (\texttt{task2\_reduce.cuh}) implements \texttt{warp\_reduce\_sum/max} with
\texttt{\_\_shfl\_xor\_sync} and \texttt{block\_reduce\_sum/max} via a
\texttt{partial[]} staging buffer and two barriers. Part~B builds
\texttt{rmsnorm\_kernel} (row-wise RMS scale with learned weights) and
\texttt{softmax\_kernel} using the online \texttt{MaxSum} merge from the guide so
each row is read once. Task~4 reuses the warp helpers for tile-wise probability sums.

\subsection{Task 3: shared-memory tiled GEMM}
Following the lecture tiled GEMM, each $16\times 16$ thread block computes a
tile of $y = \mathrm{GELU}(x W^\top + b)$ with \texttt{TILE=16}, loading
\texttt{x} and transposed \texttt{W} into \texttt{\_\_shared\_\_} buffers each
$K$-step. Padding \texttt{ws} with \texttt{[TILE][TILE+1]} avoids bank conflicts;
out-of-bounds loads use zero so every thread hits both \texttt{\_\_syncthreads()} calls.

\subsection{Task 4: tiled attention with online softmax}
Each block covers \texttt{ATTN\_WARPS=4} queries of one head; each warp streams
keys/values in tiles of 32 through dynamic shared memory (\texttt{k\_tile},
\texttt{v\_tile}, cached scaled queries). Running statistics $(m,\ell,\mathrm{acc})$
implement FlashAttention-style online softmax~\cite{dao2022flashattention} with correction
$\mathrm{corr}=\exp(m-m_{\mathrm{new}})$
when a tile raises the max. Causal masking sets scores to $-\infty$ past
\texttt{last\_key}; GQA maps \texttt{kv\_h = h / (num\_heads/num\_kv\_heads)}.
Blank~4.5 must broadcast each lane's probability with \texttt{\_\_shfl\_sync} so
every lane accumulates $\sum_r p_r v_r$ across all keys in the tile (not only its own lane).

\section{Experimental Methodology}
\label{sec:methodology}

\subsection{ASR kernels}
\textbf{Hardware:} NVIDIA RTX A6000 (sm\_86, 84 SMs; test harness reports
${\sim}768$~GB/s peak DRAM bandwidth). \textbf{Software:} \texttt{mls-asr}
conda environment (Python~3.11, PyTorch~2.10+cu128, transformers~5.0), CUDA via
\texttt{/usr/local/cuda}, repository \texttt{hw1-asr}.

\textbf{Quality metrics:} per-case \texttt{max\_err} against CPU references
(\texttt{make test1--4}) and transcription string match on \texttt{test\_audio.wav}
(\texttt{./run.sh} with all CUDA kernel groups). \textbf{Cost metrics:}
per-kernel latency/GFLOP/s from tests and end-to-end ms/token with
\texttt{--profile} shares. \textbf{Controls:} same GPU for sweeps, rebuild with
\texttt{make clean} when changing \texttt{NVCCFLAGS}, default
\texttt{-O3 -std=c++17 -lineinfo} except the \texttt{--use\_fast\_math} experiment.

\subsection{RAG evidence-selection study}
\textbf{System:} practice/dev audio questions $\rightarrow$ ASR transcript
$\rightarrow$ query encoding $\rightarrow$ retrieval over 24{,}000 enc\_a document
vectors $\rightarrow$ evidence selection/context construction $\rightarrow$ ELM
answer generation (\texttt{utter-project/EuroLLM-22B-Instruct-2512}, temperature
0.1). Retrieval uses exact cosine similarity with saved transcripts to isolate
downstream changes.

\textbf{Baseline and evaluation:} on all 350 practice queries we report
Recall@5/10/20/100, MRR@10 and NDCG@10 against public qrels. A small 5-query
automatic answer-judge sample is treated as diagnostic only.

\textbf{Held-out dev comparison (frozen method):} on 125 dev queries we compare
five configurations using the same transcripts, corpus, model and prompting:
\texttt{dense5} (top-5), \texttt{dense20} (top-20), \texttt{rerank5} (cross-encoder
reranked top-5~\cite{nogueira2019passage}), \texttt{dense\_budget} (dense fill to
${\sim}$4096 input tokens), and \texttt{rerank\_budget} (reranked fill to
${\sim}$4096). We report automatic
judge scores (correctness, grounding, completeness in $[0,4]$ and normalised
answer score in $[0,1]$), selected-id qrel recall, average selected documents and
input tokens, truncation counts, and mean reranking/generation times. This is a
single-run confirmation on dev; no tuning was performed on dev after scores were
observed. Automatic judge and qrels are imperfect proxies, and remote generation
latency is load-dependent, so timing/cost results are interpreted as partial.

\section{Results and Analysis}
\label{sec:results}

\subsection{RAG baseline and retrieval diagnostics}
We ran the supplied RAG pipeline on an ICF RTX A6000 with two allocated CPUs and
24~GB host memory (job 3713067, 11 November 2026; code revision \texttt{c0b7a7f}):
all custom CUDA ASR groups, CPU enc\_a query encoding, exact cosine retrieval over
24{,}000 normalised 384-dimensional enc\_a vectors~\cite{reimers2019sbert}, and
top-5 evidence passed to ELM EuroLLM-22B-Instruct-2512~\cite{martins2024eurollm}
(temperature 0.1).

On all 350 practice queries (public qrels, top-100 retrieved per query),
Recall@5/10/20/100 was 0.421/0.480/0.579/0.770, MRR@10 was 0.489, and NDCG@10 was
0.420. We saved ASR transcripts for controlled downstream comparisons. A 5-query
automatic answer-judge sample (1.43\% of practice) gave mean
correctness/grounding/completeness 1.6/2.0/1.2 out of 4 (normalised 0.40); one
judge error on the first query led us to treat these scores as diagnostic only.
The gap between Recall@100 and Recall@5 motivated evidence selection. In an earlier
30-query practice pilot, top-20 raised selected-ID recall (${\sim}0.29$ to
${\sim}0.51$) at about $3.7\times$ input tokens, while CPU cross-encoder reranking
did not beat dense top-5; bootstrap intervals crossed zero. We therefore froze the
method and used dev only as held-out confirmation.

\subsection{RAG held-out dev validation: evidence selection and context budget}
We evaluated five frozen configurations on 125 dev queries as a held-out
confirmation (no tuning on dev after scores). Table~\ref{tab:rag-dev} summarises
quality, evidence coverage, and context/cost indicators.

\begin{table}[t]
\centering
\small
\caption{Held-out dev validation (125 queries, single run). Quality scores are
automatic judge means; recall is selected-id qrel recall. Rerank/generation times
are partial (remote generation is load-dependent).}
\label{tab:rag-dev}
\begin{tabular}{@{}lrrrrrr@{}}
\toprule
Arm & Ans. & Recall & Docs & Tok. & Trunc. & Gen. ms \\
\midrule
\texttt{dense5}          & 0.331 & 0.379 & 5.0  & 1209 & 0   & 4622 \\
\texttt{dense20}         & 0.331 & 0.580 & 20.0 & 4948 & 0   & 5479 \\
\texttt{rerank5}         & 0.309 & 0.333 & 5.0  & 1470 & 0   & 4469 \\
\texttt{dense\_budget}   & 0.328 & 0.566 & 18.2 & 4096 & 118 & 6738 \\
\texttt{rerank\_budget}  & 0.283 & 0.476 & 15.3 & 4096 & 120 & 6450 \\
\bottomrule
\end{tabular}
\end{table}

\textbf{Findings.} (1) \emph{Higher evidence coverage did not guarantee higher
answer quality.} \texttt{dense20} and \texttt{dense\_budget} increased qrel recall
substantially over \texttt{dense5} (0.580/0.566 vs.\ 0.379) but achieved similar
normalised answer scores (0.331/0.328 vs.\ 0.331). (2) \emph{Reranking was not
beneficial in this setting.} \texttt{rerank5} was slightly below baseline, and
\texttt{rerank\_budget} had the lowest score (0.283) despite moderate recall;
reranking added ${\sim}262$~ms/query on average. (3) \emph{Budget filling changed
the quality--coverage--cost trade-off.} Budget arms used near-cap contexts with
frequent truncation, increasing recall and context/timing cost but not reliably
improving judged answers.

\textbf{Interpretation.} The dev results suggest that answer quality depends not
only on retrieving more annotated evidence but also on which evidence is selected,
how it is composed under a token limit, and whether the model uses it faithfully.
In some cases, extra context can include useful information missed by top-5; in
others, larger or reranked contexts appear to dilute or displace useful evidence.
For example, on one dev query about government bonds, budget filling included a
later useful document and raised the judged score from 0 to 0.583, whereas the
reranked budget answer remained at 0; on another query about telecommuting tax,
top-5 already scored 1.0, but larger/reranked contexts scored 0 despite sharing
core documents. These single cases are illustrative only. Because this is a single
held-out run with automatic judging and qrel-based recall, we treat the overall
pattern as observed trade-offs rather than causal or statistically significant
effects.

\subsection{Task 1 correctness and end-to-end check}
All four \texttt{make test1} cases pass with \texttt{max\_err} $\leq 2\times 10^{-7}$
(default flags). \texttt{./run.sh --cuda task1} yields 100\% transcription accuracy
(\emph{Concord returned to its place amidst the tents.}) with \texttt{rope=cuda}
and remaining groups on PyTorch.

\subsection{Selected ASR question answer (Task 1)}
\textbf{Selected question (Task 1, encoder memory bandwidth).}
The harness counts bytes moved as
$2 \cdot (\text{q\_rows} + \text{k\_rows}) \cdot \text{rotary\_dim} \cdot 4$
(read $+$ write over the rotated prefix). For the audio-encoder case
($B{=}1$, $\mathrm{seq}{=}1500$, $H_q{=}H_k{=}20$, $\text{rotary\_dim}{=}32$),
$\text{q\_rows}{=}\text{k\_rows}{=}30{,}000$, so
$2 \cdot 60{,}000 \cdot 32 \cdot 4 = 15.36 \times 10^6$~bytes (${\sim}14.6$~MiB)
per call. At ${\sim}0.026$~ms this is ${\sim}601$~GB/s, about \textbf{78\%} of the
quoted 768~GB/s peak---a roofline-style~\cite{williams2009roofline} check
consistent with a memory-bound kernel that also executes trigonometry and indexing.

\subsection{Additional Task 1 analysis (optional)}
\textbf{Additional question (precomputed $\cos/\sin$ table).}
A table of shape $[\mathrm{seq\_len}, h]$ for $\cos$ and $\sin$ adds
$2 \cdot \mathrm{seq\_len} \cdot h \cdot 4$~bytes read per pair while removing
\texttt{cosf}/\texttt{sinf} FLOPs. For the encoder case ($\mathrm{seq\_len}{=}1500$,
$h{=}16$) this is $192$~KiB total but \textbf{+8~bytes/pair} on the critical path.
At ${\sim}78\%$ of peak already, extra reads likely slow the kernel here;
precomputation is more attractive when trigonometry dominates and memory headroom
is larger.

\textbf{Additional question (single launch for $q$ and $k$).}
The model invokes RoPE \textbf{396 times} per transcription. Fusing $q$ and $k$
halves kernel-launch count, avoids an extra global synchronisation point, and
reuses one grid configuration, at the cost of slightly more complex indexing
(blank~1.2). This is favourable because launch overhead accumulates across many
invocations.

\subsection{Optional Task 1 exploration}
Table~\ref{tab:rope-threads} sweeps \texttt{ROPE\_THREADS} on the encoder case
(default \texttt{NVCCFLAGS}). Throughput is flat near 76--79\% of peak for
block sizes 64--512; \texttt{1024} drops slightly (fewer blocks, less latency hiding).
\texttt{32} is slowest (about 76\% of peak): many small blocks increase launch and
scheduling overhead relative to useful memory work.

\begin{table}[t]
\centering
\caption{RoPE encoder-case bandwidth vs.\ \texttt{ROPE\_THREADS} (RTX A6000).}
\label{tab:rope-threads}
\begin{tabular}{@{}rrl@{}}
\toprule
\texttt{ROPE\_THREADS} & Time (ms) & Effective bandwidth \\
\midrule
32  & 0.026 & 581 GB/s, 0.76 of peak \\
64  & 0.026 & 598 GB/s, 0.78 of peak \\
128 & 0.025 & 606 GB/s, 0.79 of peak \\
256 & 0.026 & 601 GB/s, 0.78 of peak \\
512 & 0.026 & 599 GB/s, 0.78 of peak \\
1024 & 0.026 & 591 GB/s, 0.77 of peak \\
\bottomrule
\end{tabular}
\end{table}

\textbf{Tiny cases (${\sim}0.002$~ms).}
Micro-benchmarks move only hundreds of floats; time is dominated by \textbf{fixed
launch latency} and synchronisation, not by bytes transferred, so reported GB/s
is near 0\% of peak even when correct.

\textbf{\texttt{--use\_fast\_math}.}
Rebuilding with \texttt{NVCCFLAGS} including \texttt{--use\_fast\_math} leaves tiny and
decoder cases passing but \textbf{fails the audio-encoder case}
(\texttt{max\_err}${\approx}1.9\times 10^{-4}$ vs.\ tolerance $10^{-5}$).
Encoder tokens reach $\mathrm{pos}\uparrow 1499$ with $\mathrm{inv\_freq}[0]{=}1$,
so $\angle$ is large; hardware-style \texttt{cosf}/\texttt{sinf} approximations
drift from the CPU reference. Decoder angles stay smaller ($\mathrm{seq}{=}64$),
so fast math still passes. Default flags were restored after the experiment
(\texttt{make clean} then \texttt{make}).

\subsection{Tasks 2--4: tests and full ASR run}
\texttt{make test2} through \texttt{make test4} all report \textbf{ALL PASS} on the
RTX A6000 (encoder attention ${\sim}8.6$~ms, ${\sim}1338$~GFLOP/s in the harness).
\texttt{./run.sh} with every group on CUDA transcribes correctly (${\sim}1961$~ms per
utterance vs.\ ${\sim}475$~ms with only Task~1 on CUDA---GEMM dominates). Profile:
\texttt{linear} ${\sim}82.5\%$ of GPU time, \texttt{attention} ${\sim}15.2\%$,
remaining norms/activations under $2\%$ each.

\subsection{Selected ASR question answers (Tasks 2--4)}
\textbf{Selected question (Task 2, reductions/normalisation/softmax).}
Reductions combine per-thread partials without atomics; RMSNorm requires one
block-wide sum of squares per row; and online softmax avoids multiple passes by
maintaining running max/sum statistics. Warp shuffles are preferred inside a warp
over shared-memory trees because they avoid shared-memory round-trips for the
intra-warp stage.

\textbf{Selected question (Task 3, tiled GEMM).}
Encoder \texttt{fc1} alone performs $1500\times5120\times1280$ MACs; tiled GEMM
reduces global memory traffic by reusing each loaded \texttt{x}/\texttt{W} element
across \texttt{TILE} threads via shared memory. The implementation reaches
hundreds of GFLOP/s on large cases but remains below cuBLAS because it does not
use tensor cores or highly tuned libraries such as CUTLASS.

\textbf{Selected question (Task 4, tiled attention/online softmax).}
A dense encoder score tensor would require
$20 \times 1500 \times 1500 \times 4 \approx 172$~MiB per layer, whereas tiled
attention streams 32-key tiles in ${\sim}48$~KiB shared memory per block. All warps
must follow the same tile iterations due to \texttt{\_\_syncthreads()} (so per-warp
early exit is unsafe). Online-softmax correction uses $\mathrm{corr}=1$ when a tile
does not raise $m$, and guards the uninitialised-$m$ case to avoid NaNs. With four
queries per block each key tile is amortised across four queries; larger
\texttt{ATTN\_WARPS} improves reuse but increases shared-memory/register pressure.
A key correctness detail was broadcasting each lane's probability with
\texttt{\_\_shfl\_sync} so every lane accumulates $\sum_r p_r v_r$ over the tile.

\section{Discussion and Conclusion}
\label{sec:conclusion}

\textbf{ASR.} All four kernel groups pass unit tests and reproduce the reference
transcription with CUDA enabled. Measurements and profiling indicate RoPE is
memory-bound, GEMM dominates end-to-end latency, and attention is the next largest
cost but benefits from tiled/online-softmax execution. Key implementation lessons
include correct index arithmetic for fused RoPE, warp/block synchronisation for
reductions and softmax, and cross-lane probability broadcasting in Task~4; fast
math is unsafe for RoPE at large encoder angles.

\textbf{RAG.} The practice baseline showed a large Recall@100--Recall@5 gap, but the
held-out dev comparison indicates that increasing or reranking evidence under a
fixed budget does not automatically improve judged answer quality. In this setting,
\texttt{dense20} and \texttt{dense\_budget} improved evidence coverage over
\texttt{dense5} without improving answer score, while \texttt{rerank\_budget}
performed worst. This supports a quality--coverage--cost trade-off view: effective
evidence selection depends on relevance, context composition, truncation, and
faithful generation, not coverage alone.

\textbf{Limitations.} ASR correctness is evidenced by provided tests and a reference
audio case, not exhaustive dataset-level transcription evaluation. RAG conclusions
are based on one held-out dev run (125 queries), automatic judge/qrel proxies, and
partial cost measurements (remote generation latency varies; no fully matched
end-to-end ASR+RAG timing). We therefore avoid causal or statistically significant
claims and report dev as confirmation of frozen methods only.

\textbf{Overall.} The combined system work shows that local kernel optimisations and
downstream RAG design choices must be evaluated with both quality and cost in mind:
correct low-level implementations enable the pipeline, but end-task gains depend on
how evidence is selected and used under realistic constraints.

% References are excluded from the main-body page limit.
\bibliography{references}
\bibliographystyle{iclr2027_conference}

% Put your Appendix here if needed. Comment out if you don't need it.
\appendix
\section{ASR reproduction commands}
\label{app:repro}

\begin{verbatim}
conda activate mls-asr
export PATH=/usr/local/cuda/bin:$PATH
cd edin-mls-26-autumn/hw1-asr
make test          # or make test1 .. test4
./run.sh           # all CUDA kernels
./run.sh --profile
\end{verbatim}

\section{RAG reproduction record}
Baseline: \texttt{rag/experiments/baseline-3713067.md} and
\texttt{rag/results/icf-3713067/}. Practice pilot:
\texttt{rag/experiments/selection-pilot-results.md} and
\texttt{rag/experiments/selection-pilot-metrics.json}. Held-out dev validation
(frozen methods, no post-hoc tuning): \texttt{rag/results/dev-validation/}
(\texttt{comparison.json}, \texttt{answers/}, \texttt{plan/}, logs). Procedure:
\texttt{rag/experiments/DEV\_VALIDATION.md}. ASR logs:
\texttt{hw1-asr/results/cuda-submit-logs/} when present.

% Keep this line unchanged.
\newpage

% Academic integrity is an underlying principle of research and academic practice.
% All submitted work is expected to be your own, including any material you have
% generated yourself with the assistance of AI tools such as ELM, ChatGPT,
% Bing Chat, Claude, Google Gemini, or any other automated assistance.
% For all such material you must make prominent explicit acknowledgement and
% explain clearly what you have done in this section.
% This includes using AI to directly generate text, code, figures or similar
% or to complete other tasks. In all cases, you remain entirely responsible
% for the material you produce.
\section*{AI declaration}
AI coding assistants (including Cursor/ELM) were used to help run and organise
GPU/RAG experiments, derive bandwidth and cost summaries, analyse saved
dev-validation results, and draft/edit LaTeX text for this report. Reported numbers
were reviewed against repository artifacts, kernel changes were checked with
available tests and runs, and the final wording and structure were edited by the
authors. The group remains responsible for the submitted work.

% Please put the information of your group here without changing the format.
% This page will be deleted when the report is distributed for reviewing.
\section*{Group information}
Group ID: XXX

Information of group member A:
\begin{itemize}
    \item Name: jacksu
    \item UUN: s2510156
\end{itemize}

Information of group member B:
\begin{itemize}
    \item Name: TBD (partner)
    \item UUN: TBD (partner)
\end{itemize}

% Keep this COMMENTED OUT, unless there is more than one member and you feel like
% it is necessary to point out the contribution. Please notice that members from
% the same group technically receive the same score for the report, regardless of
% the contribution.
% \section*{Contribution}

\end{document}
